\chapter{Implementación}
\label{cap:implementacion}
Este capítulo describe cómo se concretan los componentes del diseño en la
candidata Neuma 0.12.20, código 35, cuya base quedó identificada por la etiqueta
\path{v0.12.20-preview.1}, y las correcciones locales de 0.12.21--0.12.23
del 12 de noviembre de 2026, la comparación visual 0.12.24 y el primer
incremento de acabado Music Garden del 16 de noviembre. La candidata local
más reciente es 0.12.27, código 42; estas ampliaciones no implican una nueva
publicación de la etiqueta.
Organizo la explicación por funcionamiento del
sistema, conservando los incrementos como referencias de trazabilidad. Las
herramientas se presentan en el capítulo~\ref{cap:herramientas} y los resultados
de las verificaciones, en el capítulo~\ref{cap:pruebas}.

\section{Organización del código}
El módulo de aplicación reside en \path{Plataforma/app/}. Las fuentes Kotlin
se organizan bajo \path{ar/edu/unmdp/entrenamiento/}: \texttt{ui} reúne las
pantallas; \texttt{voice}, \texttt{instruments}, \texttt{rhythm} y
\texttt{defender}, los componentes de las actividades; \texttt{dda}, el motor
difuso; y \texttt{storage}, la coordinación de datos. Algunos ViewModels se
ubican junto a su actividad y otros en \texttt{ui}; la responsabilidad se
reconoce por el contrato, sin presuponer una distribución uniforme de archivos.

Los recursos empaquetados incluyen muestras, fuentes y grabado musical en
\path{src/main/assets/}, además de recursos Android en \path{src/main/res/}.
Las pruebas locales están en \path{src/test/} y las instrumentadas, en
\path{src/androidTest/}. Los scripts de \path{Plataforma/tools/} preparan recursos
y reúnen verificaciones. Esta estructura materializa las responsabilidades
de la sección~\ref{sec:arquitectura-neuma}; no representa varios módulos
desplegados por separado.

\section{Captura y análisis de la respuesta vocal}
\label{sec:implementacion-audio}
El módulo de captura y análisis se apoya en \path{voice/AndroidVoiceCapture.kt} y
\path{voice/YinEstimator.kt}, bajo el paquete de la aplicación, y en el laboratorio
\path{voice/PitchLab.kt}. No incorpora nuevas mediciones ni una calibración de
voz real. Los fundamentos se presentan en la sección~\ref{sec:audio-estimacion}.

La captura utiliza \texttt{AudioRecord} con un canal y muestras PCM de 16 bits,
lecturas no bloqueantes y procesamiento fuera del hilo de interfaz. La API devuelve
la cantidad leída o un error; el mínimo de búfer no garantiza continuidad bajo carga
\parencite{androidAudioRecord}. El código acumula las lecturas hasta completar el
bloque nuevo y arma la ventana del estimador. La frecuencia de muestreo y el formato
deben ser coherentes en ambos componentes.

\label{sec:ventanas-audio}
Para analizar un flujo continuo se toma un bloque finito de muestras y se repite
el cálculo al incorporar datos nuevos. Se denomina aquí ventana de análisis al
bloque de \(N\) muestras utilizado para una estimación y salto a las \(H\) muestras
que separan el comienzo de dos bloques consecutivos. En este contexto, ventana
designa el segmento analizado, sin suponer que se aplique una función de ponderación.
Las duraciones correspondientes se obtienen directamente de la frecuencia de
muestreo:
\begin{equation}
  T_{\mathrm{ventana}}=\frac{N}{f_s},\qquad
  T_{\mathrm{salto}}=\frac{H}{f_s}.
  \label{eq:ventana-salto}
\end{equation}
La configuración \texttt{yin-v1} utiliza \(f_s=44\,100\) Hz, \(N=4096\) y
\(H=1024\): aproximadamente 92,9 ms de muestras por ventana y 23,2 ms por salto.
Son duraciones calculadas desde la configuración, no mediciones de latencia.
Si \(H<N\), los bloques comparten muestras. Una mayor frecuencia de actualización
no implica que cada resultado proceda de un segmento independiente ni que el
sistema responda con una demora igual al salto.

Esta distinción permite formular el compromiso temporal del proyecto. Para
buscar un período deben estar disponibles las muestras que exige su comparación;
las frecuencias más graves corresponden a períodos más largos. Aumentar el bloque
aporta una historia temporal más extensa, pero también puede reunir parte de una
transición entre notas. La elección de tamaño, salto y rango debe comprobarse
con las señales y tareas previstas, conservando su configuración en los registros.


\subsection{Algoritmo basado en YIN}
\label{sec:estimacion-yin}
YIN estima la frecuencia fundamental mediante comparaciones temporales de la
señal consigo misma. Su función de diferencias acumula discrepancias cuadráticas
para un desplazamiento \(\tau\); la normalización por media acumulada facilita
la búsqueda del período \parencite[pp.~1919--1920]{decheveigne2002yin}.
Adaptando la notación a un bloque, para \(W\) pares de muestras disponibles:
\begin{equation}
  d(\tau)=\sum_{j=0}^{W-1}\bigl(x[j]-x[j+\tau]\bigr)^2,
  \qquad
  d'(\tau)=\frac{d(\tau)}{\frac{1}{\tau}\sum_{k=1}^{\tau}d(k)}.
  \label{eq:yin-diferencias}
\end{equation}
La segunda expresión corresponde a \(\tau>0\) y denominador positivo; para
\(\tau=0\) se define \(d'(0)=1\). El método busca un mínimo normalizado bajo un
umbral y refina el período mediante interpolación parabólica. A partir del
desplazamiento estimado \(\widehat{\tau}\), expresado en muestras, se obtiene
\(\widehat{f}_0=f_s/\widehat{\tau}\)
\parencite[pp.~1920--1921]{decheveigne2002yin}.

La versión del prototipo debe describirse como una implementación basada en YIN.
Incluye diferencias, normalización, umbral e interpolación, pero adopta decisiones
propias: rechaza el bloque si no encuentra un candidato bajo el umbral, interpola
la función normalizada y no incorpora la búsqueda adicional del mejor candidato
en un entorno temporal. El artículo contempla un mínimo global alternativo,
refinamiento sobre la diferencia sin normalizar y una etapa de búsqueda temporal
\parencite[pp.~1920--1921]{decheveigne2002yin}. Por ello, sus resultados publicados
no se atribuyen automáticamente al código de Neuma.


\subsection{Admisión de estimaciones y evaluación musical}
\label{sec:admisibilidad-tono}
Una estimación de frecuencia es un dato intermedio. Para evaluar una respuesta
vocal, el proyecto necesita decidir primero si dispone de señal admisible y,
después, cómo compara esa estimación con el objetivo musical. La ausencia de tono
aceptado debe conservar un significado distinto de una nota válida fuera de la
tolerancia de afinación.

El prototipo utiliza el nivel RMS del bloque, calculado tras retirar su media,
como una condición de admisión. Para muestras normalizadas, la expresión es
\begin{equation}
  \overline{x}=\frac{1}{N}\sum_{n=0}^{N-1}x[n],\qquad
  x_{\mathrm{RMS}}=\sqrt{\frac{1}{N}\sum_{n=0}^{N-1}
  \bigl(x[n]-\overline{x}\bigr)^2}.
  \label{eq:rms-bloque}
\end{equation}
Esta magnitud describe el nivel numérico de la señal analizada; sin calibración
del sistema de captura no representa una medición de presión sonora. Superar su
umbral tampoco identifica por sí mismo una voz ni garantiza periodicidad.

Además del nivel, la implementación comprueba la validez de las muestras, la
existencia de un candidato y su pertenencia al rango configurado. El indicador
derivado del mínimo normalizado se interpreta como un descriptor del cálculo,
sin presentarlo como probabilidad calibrada de que la nota sea correcta. El
criterio musical se formula luego mediante distancias como los cents de la
sección~\ref{sec:afinacion-cents} y, cuando corresponda, condiciones temporales
de estabilidad. Los umbrales de señal, de estimación y de aprobación musical
cumplen funciones diferentes y deben evaluarse por separado.

\section{Integración de intervalos vocales}
\label{sec:implementacion-intervalos}
\subsection{Recorrido de captura y evaluación}
\texttt{VocalIntervalViewModel} conecta referencia sonora, capturador, sesión
y registro. La imitación reproduce dos notas y recibe dos capturas nuevas:
primero el ancla cantada y después el objetivo, mediante acciones explícitas.
Ancla dada reproduce una nota y captura solo el objetivo; producción libre
recibe ambas notas sin reproducción previa. En los tres modos la sesión
conserva la procedencia del ancla y evalúa la relación que corresponde.

Con ancla cantada, el intervalo se calcula como
\(1200\log_2(f_2/f_1)\) y se compara con el objetivo en semitonos. Trasladar
ambas notas al mismo tiempo conserva el intervalo, mientras que un desplazamiento
de octava en una sola conserva su error. Con ancla dada, la comparación utiliza
la frecuencia de la referencia reproducida, sin inventar mediciones de una
primera nota no cantada.

La sesión adquiere notas mediante sostenido y mediana en escala logarítmica.
La configuración experimental utiliza 500 ms de sostenido, huecos máximos de
100 ms y una amplitud de 35 cents para la estabilidad. Una segunda nota estable
puede cerrar con error de afinación; silencio o discontinuidad reinician la
adquisición. La tolerancia de aprobación es un parámetro distinto de la
admisibilidad de señal descrita en la sección~\ref{sec:admisibilidad-tono}.

El modelo serializa acciones y estimaciones. Una generación invalida las
respuestas tardías de operaciones canceladas; la siguiente operación espera
la liberación del recurso anterior. Cada etapa abre una captura nueva para
evitar arrastrar muestras del ancla al objetivo. Conceder permiso no inicia
por sí solo el micrófono. Cancelación, salida o segundo plano interrumpen sin
reanudar automáticamente, y un temporizador limita también la espera en silencio.

\subsection{Adaptación, series y registro}
La progresión organiza series de cinco rondas del mismo intervalo, dirección
y ayuda. La práctica manual permite ensayar contenidos sin conceder avance.
El DDA considera el error absoluto medio y la dispersión del sostenido final,
mediante el motor compartido. La tolerancia experimental comienza en 50 cents
y cambia en pasos de hasta 5, dentro de 30--80 cents; el ajuste se utiliza
en la ronda siguiente, sin modificar la que está siendo evaluada.

Dos series altas elegibles permiten registrar el hito del contenido. Imitación,
ancla dada y producción libre conservan tolerancia, racha y evidencia propias.
Un hito puede habilitar acceso a la ayuda siguiente sin acreditar respuestas
en ella. La selección entre las tres ayudas sigue siendo explícita; RF-16
representa la gobernanza automática todavía por implementar.

\texttt{VocalIntervalHistorySession} asigna identidad al comenzar la captura y
registra cierres completos e interrupciones. Escuchar sin cantar no crea una
ronda. El registro conserva referencia, ayuda, configuración de análisis,
dispositivo y serie derivada, separado del laboratorio técnico de voz.
En producción libre, las referencias MIDI son nulas y las escuchas, cero.
La decisión y los estados de progreso se preparan con el guardado del mismo
registro; el reintento conserva su identidad. Los perfiles anteriores cargan
los mapas nuevos vacíos, sin transferir desempeño de otras ayudas.

\section{Reconocimiento de instrumentos}
\label{sec:implementacion-instrumentos}
\texttt{InstrumentSession} prepara una melodía al crear cada pregunta y la
mantiene en las reescuchas. El reproductor compone la frase con muestras locales
del instrumento y conserva notas, duraciones y versión del generador para el
resultado. Las opciones y el instrumento objetivo pertenecen al catálogo del
nivel; la misma gramática de melodías se utiliza entre instrumentos para
reducir diferencias de tarea ajenas al reconocimiento tímbrico.

La política intrapartida utiliza aciertos y autonomía de escucha, sin penalizar
el tiempo de respuesta. Desde la tercera respuesta examina una ventana de
tres preguntas; puede cambiar una escucha para la siguiente pregunta, entre
una y tres, y deja dos preguntas entre cambios. El ejemplo ya presentado
conserva su ayuda. Un acierto con reescuchas sigue siendo correcto, aunque
aporta información distinta sobre autonomía.

El modo manual conserva la ayuda fija de práctica; el automático aplica la
política y aporta evidencia a la progresión del nivel. El resultado reúne
respuestas, melodías, escuchas y trazas de decisión. Los registros históricos
de nota aislada conservan su identificación y no reciben una melodía ficticia.
El inspector y los escenarios controlados permiten revisar esas políticas
con datos simulados separados del perfil.

Desde 0.12.21 se permite responder durante la melodía, una vez notificado el
inicio real del audio. La elección cancela la reproducción y conserva una sola
respuesta y la cantidad de escuchas iniciadas, sin penalización por rapidez.

\section{Generación y ejecución de Rhythm Tap}
\label{sec:implementacion-ritmo}
\subsection{Partitura, entrada y revisión}
El modelo rítmico conserva duraciones racionales, figuras, silencios, ligaduras
y procedencia de categorías. El generador comprueba selección y compás antes
de buscar un ejercicio, con cantidad de candidatos y tiempo de búsqueda
acotados. El agotamiento de búsqueda se comunica como un fallo recuperable,
distinto de una selección que no cumple las condiciones de factibilidad.

La partitura se presenta antes de iniciar. Su modelo alimenta notación, ataques
esperados y evaluación, de modo que la continuación de una ligadura no produzca
un ataque adicional. El grabado local se integra en \texttt{RhythmScoreView};
el desplazamiento lateral y el marcador siguen la ejecución, con una alternativa
de ventanas estáticas al reducir movimiento.

\texttt{RhythmViewModel} relaciona la marca temporal del evento táctil con el
reloj de ejecución. La evaluación distingue ataques adelantados o atrasados,
omisiones y golpes extra. Repetir conserva la partitura como práctica; generar
otro identificador no convierte el mismo contenido en evidencia nueva de avance.

\subsection{Progresión y revisiones del resultado}
\label{sec:implementacion-ritmo-r6}
En progresión, cada categoría conserva tempo recomendado y evidencia propios.
Al mezclarlas se utiliza el menor tempo recomendado, fijo durante la ejecución.
Dos resultados altos elegibles permiten aumentar cinco BPM para un ejercicio
posterior y acreditar el hito de la categoría. Inicio a 60 BPM, techo de
120 BPM y prerrequisitos son parámetros experimentales.

La inferencia utiliza precisión de ataques y silencios, y control de golpes
extra atribuibles, considerando las apariciones objetivo de cada categoría.
Una huella del contenido musical independiente de semilla y UUID identifica
las repeticiones. La categoría más rápida no gana tempo por ejecutarse al
tempo recomendado de otra más lenta.

\texttt{prepareRhythmCommit} prepara resultado, decisión y progreso dentro del
guardado coordinado. Si se corrige el último intento, recalcula desde su estado
anterior conservando avances adquiridos. Cuando ya se confirmó uno posterior,
actualiza la evaluación musical sin cambiar la decisión utilizada para continuar;
la revisión evaluada identifica esa evidencia. Esta política evita repetir
retroactivamente cambios de dificultad sobre ejercicios ya ejecutados.

\subsection{Conteo y alineación previa del toque}
En 0.12.22 el conteo visual utiliza el tempo de la partitura; se corrigió una
expresión que conservaba 60 BPM. La entrada recibe los eventos táctiles nativos
con su marca original. El metrónomo mantiene una línea PCM y una preparación
silenciosa de 750 ms antes del primer pulso; el sonido opcional de cada tap
se precarga en \texttt{SoundPool}, evitando pausar y recargar una pista por golpe.
Estas elecciones requieren comprobación acústica en el dispositivo.

La calibración es una actividad separada: cuatro pulsos de escucha y doce
de toma a 80 BPM. Propone la mediana del desfase y comprueba la regularidad
de los contactos antes de ofrecer su aplicación explícita. El ajuste se resta
del tiempo relativo al evaluar, sin cambiar las ventanas ni aprender de
partidas puntuadas. El resultado conserva las muestras y el ajuste junto con
tiempos originales y corregidos. Esta alineación incluye sincronización personal
y respuesta del equipo; no permite aislar latencia física. Se conserva en memoria
y debe repetirse al cambiar la salida de audio o iniciar un nuevo proceso.

\section{Lectura musical en Defender}
\label{sec:implementacion-defender}
\subsection{Tiempo de entrada y cierre de partida}
\label{sec:implementacion-defender-d3}
\texttt{DefenderSession} evalúa la marca temporal original de la pulsación frente
al plazo de llegada del objetivo, independientemente del instante en que se
procesa o anima. Los eventos permiten corregir una llegada cuando la entrada
ocurrió antes; esa corrección puede deshacer el cierre por la tercera vida.
Por ello, el resumen inicial se guarda como provisional, sin conceder progreso.
Al cerrar su entrada mediante navegación se sella una revisión del mismo
intento y se confirma su avance. Abandonar después de una corrección conserva
la interrupción correspondiente.

El DDA observa los últimos cinco objetivos resueltos, incluidas las llegadas,
y combina supervivencia, precisión y fracción del viaje consumida. Sus ajustes
modifican objetivos futuros; los visibles conservan condiciones asignadas.
La progresión exige partida completa y evidencia pertinente al contenido,
separada de la práctica manual.

La configuración local de 0.12.22 permite viaje mínimo de 3 s, solicitado por
el autor tras probar el prototipo. La progresión conserva inicio en 10 s,
paso de 1 s y cambios solo para enemigos futuros. En práctica se puede elegir
intervalo de aparición fijo entre 1 y 6 s, en pasos de 0,5 s, o conservar la
curva automática existente de 4 a 2,5 s según tiempo activo. Ese intervalo
manual no se traslada a progresión. Los rangos siguen siendo experimentales.

\subsection{Grafía, respuesta y sonido de acordes}
\label{sec:implementacion-defender-d4}
El modelo conserva raíz escrita, calidad y voces con octava. La construcción
determina letras de terceras, quintas y séptimas antes de asignar sus alteraciones.
El catálogo excluye dobles alteraciones y filtra las disposiciones que exceden
el rango escrito admitido. Notación, solución y sonido se derivan del mismo
contenido.

La respuesta por nombre utiliza raíz, alteración y calidad, sin octava: dos
grafías enarmónicas pueden compartir alturas y exigir respuestas diferentes.
Las voces del acorde se reproducen al acertar mediante síntesis simultánea
de tres o cuatro sinusoides. La notación alterna cabezas en segundas escritas
y dispone alteraciones en columnas. La respuesta mediante piano de RF-07
permanece pendiente; esas reglas por nombre no sustituyen la evaluación futura
del conjunto exacto de teclas.

\subsection{Inversiones y progreso por contenido}
\label{sec:implementacion-defender-d5}
Las inversiones se obtienen rotando voces elevadas una octava, sin omisiones
ni duplicaciones. La raíz conserva su identidad aunque el bajo sea otra nota.
Se comprueban grafía y límites del pentagrama antes de sortear el objetivo,
y la síntesis utiliza las alturas de la disposición mostrada.

El progreso distingue calidad, inversión y contexto de claves. La configuración
comprende 84 objetivos derivados de ocho calidades, sus inversiones y los
contextos sol, fa y mezcla. Una partida en mezcla requiere evidencia de ambas
claves. Los datos anteriores de tríadas solo aportan evidencia de fundamentales
en sol; no acreditan inversiones o claves que no se practicaron. Los mapas y
estados de progreso se conservan en el documento de Defender, con decisiones
y revisiones que permiten interpretar la compatibilidad de los registros.

\section{Ajuste del afinador complementario para cuerdas pulsadas}
La revisión local 0.12.23 responde a la observación del autor sobre Si y Mi
agudas de guitarra criolla en los tres modos. El afinador utiliza 400 ms
continuos dentro del margen y un perfil YIN con RMS mínimo 0,003; conserva
los márgenes de afinación y no modifica el análisis vocal. Son parámetros
experimentales que requieren comprobación con el instrumento real.

Libre distingue la última confirmación de una medición vigente cuando decae
la señal. En Oído, una comprobación nueva ya afinada puede avanzar sin exigir
otra pulsación; un intento fallido conserva su error inicial y exige evidencia
nueva para corregir. Los recorridos guardan la política y el análisis efectivos;
silencio y lecturas de otra cuerda no se acumulan como estabilidad.

\section{Integración de datos y estados de error}
\texttt{DemonstrationViewModel} coordina la salida de las operaciones activas
antes de preparar la exportación JSON o el reinicio. Reúne resultados confirmados y
pendientes identificados; la selección del destino y la escritura presentan
sus estados de espera, cancelación o fallo. Al reiniciar, los modelos actualizan
su estado visible después de la operación sobre el perfil.

Las actividades reutilizan esa gestión, manteniendo las políticas particulares
de cada registro. Las simulaciones sustituyen audio, tiempo o almacenamiento
en los límites correspondientes y conservan procedencia diferenciada. La
organización de documentos y sus límites de recuperación se describe en la
sección~\ref{sec:diseno-datos}; aquí su integración permite seguir el recorrido
desde una respuesta hasta su consulta o exportación.

La candidata descrita integra las cuatro actividades con los límites señalados.
Piano en Defender y selección automática de ayuda vocal siguen pendientes,
además de las verificaciones Android y físicas identificadas por versión.
Las guías R6, D3--D5 y V1--V6 conservan decisiones y evidencia incremental;
sus resultados de prueba se presentan en el capítulo~\ref{cap:pruebas}.

\section{Variantes de presentación para revisión con el director}
El 12 de noviembre de 2026 el autor solicitó conservar dos aplicaciones
instalables para comparar sus presentaciones. La revisión 0.12.24 identifica
la aplicación existente como Neuma Garden y añade Neuma Quiet Studio con
un identificador de paquete distinto. Los perfiles permanecen independientes.
La segunda variante adapta el catálogo y el recorrido de Rhythm Tap, incluida
su calibración y consulta de resultados, con superficies cálidas y un resumen
de ejecución cuyos detalles pueden desplegarse. Reutiliza los componentes de
notación y las mismas sesiones, procesamiento de audio y reglas de adaptación.
La implementación permite presentar ambas alternativas; no representa una
elección definitiva de identidad ni una evaluación de usabilidad. El alcance,
los artefactos y la recuperación de la versión previa se registran en
\path{docs/Comparacion-Garden-Quiet-Studio-v0.12.24.md}.

\subsection{Continuidad de Music Garden y acabado de Instrumentos}
El 16 de noviembre el autor informó que acordó con el director continuar con
la estética de videojuego Music Garden y el menú del libro. Posteriormente
confirmó un acabado por recorridos: preparación, resultados e historial sobre
papel, con escenarios propios durante las partidas. Este acuerdo visual no
constituye aprobación de todos los requerimientos ni de los criterios de evaluación.

La candidata 0.12.25 aplica ese primer recorrido a Instrumentos. Reutiliza el
tema, los tokens y las fuentes del sistema; incorpora un invernadero vectorial
estático, idéntico para todos los timbres y sin pistas sobre la respuesta. Los
controles admiten crecimiento del texto y desplazamiento. La selección de nivel
conserva un espacio fijo para la marca y distingue foco, selección y deshabilitado.
El historial se abre filtrado por Instrumentos y retorna a la preparación o al
resultado de origen, sin reiniciar modalidad ni nivel. Las sesiones, el audio,
la adaptación y los formatos de datos conservan su funcionamiento.

La composición nativa es una primera versión por evaluar en teléfono. Los
sonidos de interfaz y la respuesta táctil quedan pendientes de definición e
incorporación; el personaje continúa pausado. El alcance y los recursos se
registran en la especificación de acabado Music Garden y en el inventario de
recursos, sin atribuir aprobación visual final ni validación de uso.

\subsection{Terraza y organización del recorrido rítmico}
Por pedido del autor de continuar, la candidata 0.12.26 extiende el acabado a
Rhythm Tap. Una terraza vectorial estática enmarca la previa y la ejecución;
la partitura conserva su grabado desde datos y una superficie limpia. La escena
se omite en resultados y durante la ejecución con texto ampliado por encima del
130\%, para priorizar el espacio de lectura. La preparación se pliega después
de generar o terminar un ejercicio y puede volver a abrirse. Los accesos a
historial y progreso, junto con avisos de guardado y recuperación, permanecen
fuera de ese bloque plegado.

El resultado ofrece un resumen con desglose desplegable y accesos al generador
existente y a la repetición. El pad admite crecimiento con el texto y conserva
su receptor nativo durante resultados para procesar contactos entregados tarde.
Calibración, historial y progreso reutilizan papel, tema y componentes; Quiet
Studio conserva su paleta y omite la decoración Garden. El cambio se limita a
presentación y acceso a acciones existentes, sin modificar sesiones, evaluación,
audio, relojes ni persistencia. El escenario concreto y los recorridos nativos
siguen pendientes de revisión en teléfono.

\subsection{Claro de canto y presentación de intervalos vocales}
La candidata 0.12.27 incorpora una portada vocal propia y un claro de canto
estático durante referencia y captura. El editable SVG se exporta a un recurso
vectorial Android con colores comprobados contra los tokens; no contiene
instrumentos, texto ni notación musical. La escena se retira en resultados y
durante la práctica con texto ampliado por encima del 130\%. Preparación,
captura, devolución, historial y progreso reutilizan papel y paneles compartidos.
La ayuda para señal insuficiente permanece estable durante la captura, evitando
que su aparición intermitente desplace los controles.

La reorganización conserva Imitación, Ancla dada y Producción libre, así como
sus referencias, capturas, ayudas, tolerancias y datos. El permiso de micrófono
sigue separado del inicio de una captura. El contenedor compartido centraliza
áreas seguras sin cambiar la excepción de presentación de ritmo en Quiet Studio.
Se revisó la ilustración SVG en navegador; su composición nativa y los recorridos
en teléfono siguen pendientes de evaluación.
